% GENERATED FILE. Edit canonical lessons, then run npm run generate:books.
% canonical-source-hash: fnv1a64:5701c67db8c53231
% canonical-lessons: TA-C334-ampu, TA-C334-ter, TA-C334-kutiraiviran, TA-C334-muracu, TA-C334-talapati

\chapter{Arrows and Chariots}
\label{ch:ta-arrows-and-chariots}
\hlchaptermodality{334}

\begin{chapteropening}
\textbf{By the end of this chapter:} \emph{I can talk about old weapons and warriors.}

Complete the last lesson of chapter 334: I can talk about old weapons and warriors.
\end{chapteropening}

\section[\texorpdfstring{\ta{அம்பு} (ampu)}{அம்பு (ampu)}]{\ta{அம்பு} (ampu) — an arrow}

\label{lesson:TA-C334-ampu}



\begin{quote}
Before the new one: say the Tamil for a shield, then the Tamil for armour.
\end{quote}

\subsection*{You'll want to know: \ta{அம்பு}}
\textbf{\ta{அம்பு}} — \emph{ampu} — \textquotedblleft{}an arrow\textquotedblright{}.

\begin{quote}
\ta{வேடன்} \ta{ஒரு} \ta{அம்பை} \ta{எய்தான்}.
\end{quote}

\begin{quote}
\emph{vēṭaṉ oru ampai eytāṉ.}
\end{quote}

\begin{quote}
The hunter shot an arrow.
\end{quote}

\ta{அம்பு} \ta{எய்}, to shoot an arrow.

\begin{cousinweb}
\textbf{\ta{அம்பு}} (\emph{ampu}) is native Dravidian (DEDR 178): arrow.
\end{cousinweb}

\begin{grammarlens}[title={What you've built}]
One more word for telling what happened.
\end{grammarlens}

\subsection*{Your turn}
\begin{itemize}
\raggedright
  \item \emph{Say it:} \emph{ampu}
  \item \emph{Say it:} \emph{ampu}, once more
  \item \emph{Say it:} \emph{vēṭaṉ oru ampai eytāṉ.} — The hunter shot an arrow.
  \item \emph{Recall:} say the Tamil for a shield, then the Tamil for armour, then say \emph{ampu} again
\end{itemize}

\begin{tcolorbox}[breakable,colback=teal!4,colframe=teal!35!black,title={Before you move on}]
What does \ta{அம்பு} mean? (\textquotedblleft{}An arrow\textquotedblright{}.) Say it once more.
\end{tcolorbox}
\section[\texorpdfstring{\ta{தேர்} (tēr)}{தேர் (tēr)}]{\ta{தேர்} (tēr) — a chariot; a temple car}

\label{lesson:TA-C334-ter}



\begin{quote}
Before the new one: say the Tamil for armour, then the Tamil for an arrow.
\end{quote}

\subsection*{You'll want to know: \ta{தேர்}}
\textbf{\ta{தேர்}} — \emph{tēr} — \textquotedblleft{}a chariot; a temple car\textquotedblright{}.

\begin{quote}
\ta{திருவிழாவில்} \ta{பெரிய} \ta{தேர்} \ta{தெருவில்} \ta{வந்தது}.
\end{quote}

\begin{quote}
\emph{tiruviḻāvil periya tēr teruvil vantatu.}
\end{quote}

\begin{quote}
At the festival a big chariot came down the street.
\end{quote}

Today mostly a temple car pulled through the streets at festivals.

\begin{cousinweb}
\textbf{\ta{தேர்}} (\emph{tēr}) is native Dravidian (DEDR 3459): car, chariot, vehicle.
\end{cousinweb}

\begin{grammarlens}[title={What you've built}]
One more word for telling what happened.
\end{grammarlens}

\subsection*{Your turn}
\begin{itemize}
\raggedright
  \item \emph{Say it:} \emph{tēr}
  \item \emph{Say it:} \emph{tēr}, once more
  \item \emph{Say it:} \emph{tiruviḻāvil periya tēr teruvil vantatu.} — At the festival a big chariot came down the street.
  \item \emph{Recall:} say the Tamil for armour, then the Tamil for an arrow, then say \emph{tēr} again
\end{itemize}

\begin{tcolorbox}[breakable,colback=teal!4,colframe=teal!35!black,title={Before you move on}]
What does \ta{தேர்} mean? (\textquotedblleft{}A chariot; a temple car\textquotedblright{}.) Say it once more.
\end{tcolorbox}
\section[\texorpdfstring{\ta{குதிரைவீரன்} (kutiraivīraṉ)}{குதிரைவீரன் (kutiraivīraṉ)}]{\ta{குதிரைவீரன்} (kutiraivīraṉ) — a horseman}

\label{lesson:TA-C334-kutiraiviran}



\begin{quote}
Before the new one: say the Tamil for an arrow, then the Tamil for a chariot.
\end{quote}

\subsection*{You'll want to know: \ta{குதிரைவீரன்}}
\textbf{\ta{குதிரைவீரன்}} — \emph{kutiraivīraṉ} — \textquotedblleft{}a horseman\textquotedblright{}.

\begin{quote}
\ta{ஒரு} \ta{குதிரைவீரன்} \ta{செய்தியுடன்} \ta{வேகமாக} \ta{வந்தான்}.
\end{quote}

\begin{quote}
\emph{oru kutiraivīraṉ ceytiyuṭaṉ vēkamāka vantāṉ.}
\end{quote}

\begin{quote}
A horseman came quickly with the news.
\end{quote}

\ta{குதிரை}, horse + \ta{வீரன்}, warrior.

\begin{cousinweb}
\textbf{\ta{குதிரைவீரன்}} (\emph{kutiraivīraṉ}) joins native \ta{குதிரை}, \textquotedblleft{}horse\textquotedblright{} (DEDR 1711), with \ta{வீரன்}, from Sanskrit \emph{vīra}, \textquotedblleft{}hero\textquotedblright{} (CDIAL 12056).
\end{cousinweb}

\begin{grammarlens}[title={What you've built}]
One more word for telling what happened.
\end{grammarlens}

\subsection*{Your turn}
\begin{itemize}
\raggedright
  \item \emph{Say it:} \emph{kutiraivīraṉ}
  \item \emph{Say it:} \emph{kutiraivīraṉ}, once more
  \item \emph{Say it:} \emph{oru kutiraivīraṉ ceytiyuṭaṉ vēkamāka vantāṉ.} — A horseman came quickly with the news.
  \item \emph{Recall:} say the Tamil for an arrow, then the Tamil for a chariot, then say \emph{kutiraivīraṉ} again
\end{itemize}

\begin{tcolorbox}[breakable,colback=teal!4,colframe=teal!35!black,title={Before you move on}]
What does \ta{குதிரைவீரன்} mean? (\textquotedblleft{}A horseman\textquotedblright{}.) Say it once more.
\end{tcolorbox}
\section[\texorpdfstring{\ta{முரசு} (muracu)}{முரசு (muracu)}]{\ta{முரசு} (muracu) — a war drum}

\label{lesson:TA-C334-muracu}



\begin{quote}
Before the new one: say the Tamil for a chariot, then the Tamil for a horseman.
\end{quote}

\subsection*{You'll want to know: \ta{முரசு}}
\textbf{\ta{முரசு}} — \emph{muracu} — \textquotedblleft{}a war drum\textquotedblright{}.

\begin{quote}
\ta{போர்} \ta{தொடங்கியபோது} \ta{முரசு} \ta{ஒலித்தது}.
\end{quote}

\begin{quote}
\emph{pōr toṭaṅkiyapōtu muracu olittatu.}
\end{quote}

\begin{quote}
When the war began, the war drum sounded.
\end{quote}

\ta{முரசு} \ta{அறை}, to beat the drum, once a way to make news known.

\begin{cousinweb}
\textbf{\ta{முரசு}} (\emph{muracu}) is native Dravidian (DEDR 4973): a drum, a war drum, from a family of words for making a sound.
\end{cousinweb}

\begin{grammarlens}[title={What you've built}]
One more word for telling what happened.
\end{grammarlens}

\subsection*{Your turn}
\begin{itemize}
\raggedright
  \item \emph{Say it:} \emph{muracu}
  \item \emph{Say it:} \emph{muracu}, once more
  \item \emph{Say it:} \emph{pōr toṭaṅkiyapōtu muracu olittatu.} — When the war began, the war drum sounded.
  \item \emph{Recall:} say the Tamil for a chariot, then the Tamil for a horseman, then say \emph{muracu} again
\end{itemize}

\begin{tcolorbox}[breakable,colback=teal!4,colframe=teal!35!black,title={Before you move on}]
What does \ta{முரசு} mean? (\textquotedblleft{}A war drum\textquotedblright{}.) Say it once more.
\end{tcolorbox}
\section[\texorpdfstring{\ta{தளபதி} (taḷapati)}{தளபதி (taḷapati)}]{\ta{தளபதி} (taḷapati) — a commander; a general}

\label{lesson:TA-C334-talapati}



\begin{quote}
Before the new one: say the Tamil for a horseman, then the Tamil for a war drum.
\end{quote}

\subsection*{You'll want to know: \ta{தளபதி}}
\textbf{\ta{தளபதி}} — \emph{taḷapati} — \textquotedblleft{}a commander; a general\textquotedblright{}.

\begin{quote}
\ta{தளபதி} \ta{படையை} \ta{முன்னால்} \ta{நின்று} \ta{நடத்தினார்}.
\end{quote}

\begin{quote}
\emph{taḷapati paṭaiyai muṉṉāl niṉṟu naṭattiṉār.}
\end{quote}

\begin{quote}
The commander led the army from the front.
\end{quote}

Today also a general in an army.

\begin{cousinweb}
\textbf{\ta{தளபதி}} (\emph{taḷapati}) is usually traced to Sanskrit, \emph{dala}, \textquotedblleft{}a troop\textquotedblright{}, and \emph{pati}, \textquotedblleft{}lord\textquotedblright{}; Monier-Williams lists \emph{dalapati}.
\end{cousinweb}

\begin{grammarlens}[title={What you've built}]
One more word for telling what happened.
\end{grammarlens}

\subsection*{Your turn}
\begin{itemize}
\raggedright
  \item \emph{Say it:} \emph{taḷapati}
  \item \emph{Say it:} \emph{taḷapati}, once more
  \item \emph{Say it:} \emph{taḷapati paṭaiyai muṉṉāl niṉṟu naṭattiṉār.} — The commander led the army from the front.
  \item \emph{Recall:} say the Tamil for a horseman, then the Tamil for a war drum, then say \emph{taḷapati} again
\end{itemize}

\begin{tcolorbox}[breakable,colback=teal!4,colframe=teal!35!black,title={Before you move on}]
What does \ta{தளபதி} mean? (\textquotedblleft{}A commander; a general\textquotedblright{}.) Say it once more.
\end{tcolorbox}
